\section{Solyaev's iron example}\label{sup:fe}

Solyaev~\cite[section~5]{Solyaev2023} applies his scheme to iron grains with
$C_{11}=108$, $C_{12}=62$ and $C_{44}=28.3$~GPa, a grain gradient tensor
$\ell_0^2\,\tC_\gr\ot\kron$ (his eq.~(17)) and the medium $\ell^2\tC\ot\kron$ (his
eq.~(19)), lengths in units of $\ell_0$, and reports that $\ell$ grows with the grain size
towards about $1.9\,\ell_0$ (his figure~2), an increase he attributes to the misorientation of
the grains. We read eqs.~(17) and~(19) in form~II, as the Helmholtz model of his sources
(section~7 of the paper), solved his minimisation problem at the eleven grain sizes of his
figure~2, $2R/\ell_0=10$--$100$, and obtained:
\begin{itemize}
\item with the exact averaged Eshelby tensor (appendix~B of the paper):
$\ell=\NumFeExactEllMin$--$\NumFeExactEllMax\,\ell_0$,
$\nu=\NumFeExactNuMin$--$\NumFeExactNuMax$ and $\scres\le\NumFeExactResidualMax$;
\item with any of the three averaged tensors (the exact one, that of Ma et al.\ and his
printed eq.~(21)): $\ell=\NumFeAllEllMin$--$\NumFeAllEllMax\,\ell_0$ and
$\scres\le\NumFeAllResidualMax$;
\item in the closed large-grain limit, $\linf=\NumFeLimEll\,\ell_0$ and Kr\"oner's
$\nu_{\mathrm K}=\NumFeKronerNu$.
\end{itemize}
The reason is structural. For an isotropic grain, the medium with the grain's $\nu$ and
$\ell=\ell_0$ gives
\begin{equation}
\tD=\tD_\gr,\qquad \tB=\tI,\qquad \scres=0
\end{equation}
for every $R$ and any averaged Eshelby tensor, and his iron is nearly isotropic: its Zener
ratio is $\NumFeZener$, and the Voigt and Reuss shear moduli are $\NumFeMuVoigt$ and
$\NumFeMuReuss$~GPa. An assumed, strongly anisotropic grain ($C_{11}=230$, $C_{12}=135$,
$C_{44}=117$~GPa, Zener ratio $\NumFeAnisoZener$) still gives
$\ell=\NumFeAnisoEllMin$--$\NumFeAnisoEllMax\,\ell_0$ for $2R/\ell_0=10$, $20$, $50$ and
$100$.

We then tried $\NumFeMixedCount$ mixed readings: grain or medium read literally in form~I
(his index order) and the other in form~II, with each of the three averaged tensors and with
his or Ma's fourth-order average (his eq.~(20) writes $s_2I_{ijkl}$ with
$I_{ijkl}=\tfrac12(\delta_{ik}\delta_{jl}+\delta_{il}\delta_{jk})$, half the term
$s_2(\delta_{ik}\delta_{jl}+\delta_{il}\delta_{jk})$ of~\cite[eq.~(3.18)]{Ma2018}). They show a
similar trend, but none reproduces his figure~2: the root mean square difference in
$\ell/\ell_0$ is $\NumFeMixedRmsMin$--$\NumFeMixedRmsMax$, against a digitisation uncertainty
of about $0.01$. They also move $\ell$ away from $\ell_0$ for an isotropic grain, for which the
coherent reading is exact, which points to a mismatch between grain and medium rather than to
misorientation.

Two further remarks. His eq.~(22) prints the five $B$ terms without squares; read so, the
radicand is negative in parts of the domain, and we square every term. His classical
$\nu=0.348$ and $\mu=26.16$~GPa agree with each other and with his bulk modulus, but not, to the
digits printed, with Kr\"oner's estimate, $\nu_{\mathrm K}=\NumFeKronerNu$; he does not say how
they were obtained. We could not determine what was
computed for his figure~2; the implementation of section~7 is validated instead by a second,
independent implementation in components and by the limiting cases listed in section~S2.
